\BourbakiExerciseGroup

\begin{enumerate}
\def\labelenumi{\arabic{enumi})}
\tightlist
\item
  \begin{enumerate}
  \def\labelenumii{\alph{enumii})}
  \tightlist
  \item
    设 X 是拓扑空间。证明以下三个命题等价：
  \end{enumerate}
\end{enumerate}

\begin{enumerate}
\def\labelenumi{(\Alph{enumi})}
\setcounter{enumi}{16}
\tightlist
\item
  若 \(x, y\) 是 \(X\) 的任意两个不同的点，则存在 \(x\) 的一个邻域，它不含 \(y\)。
\end{enumerate}

(Q') 由单独一点组成的 X 的每个子集在 X 中闭。(Q'\,') 对每个 \(x \in X\)，x 的诸邻域之交仅由点 x 组成。

满足这些条件的空间 X 称为可达的。

\begin{enumerate}
\def\labelenumi{\alph{enumi})}
\setcounter{enumi}{1}
\tightlist
\item
  带右拓扑的有序集 X（关于这个拓扑 X 是 Kolmogoroff 空间{[}§ 1，习题 2{]}）是可达的，当且仅当 X 的任何两个不同元素都不可比较，且此时 X 上的右拓扑是离散拓扑。
\end{enumerate}

\begin{enumerate}
\def\labelenumi{\arabic{enumi})}
\setcounter{enumi}{1}
\tightlist
\item
  \begin{enumerate}
  \def\labelenumii{\alph{enumii})}
  \tightlist
  \item
    若集合 X 上的一个拓扑具有有限子基，且 X 关于这个拓扑是可达的，则 X 有限且该拓扑是离散拓扑。
  \end{enumerate}
\end{enumerate}

\begin{enumerate}
\def\labelenumi{\alph{enumi})}
\setcounter{enumi}{1}
\tightlist
\item
  若在可达空间 X 中开集的每个交都是开集，则 X 是离散的（参见§ 7，习题 6）。
\end{enumerate}

\begin{enumerate}
\def\labelenumi{\arabic{enumi})}
\setcounter{enumi}{2}
\tightlist
\item
  \begin{enumerate}
  \def\labelenumii{\alph{enumii})}
  \tightlist
  \item
    设 X 是可达空间，A 是 X 的子集，x 是 \(\overline{A}\) 中不属于 A 的一点。证明 x 的每个邻域都含有 A 的无穷多个点。
  \end{enumerate}
\end{enumerate}

\begin{enumerate}
\def\labelenumi{\alph{enumi})}
\setcounter{enumi}{1}
\item
  设 X 是可达空间，A 是 X 的子集。证明由这样的点 \(x \in \overline{A}\) 组成的集在 X 中闭：x 的每个邻域都含有 A 的异于 x 的点。
\item
  设 X 是可达空间。证明 X 的任何子集 A 的一切邻域之交等于 A。
\end{enumerate}

\begin{enumerate}
\def\labelenumi{\arabic{enumi})}
\setcounter{enumi}{3}
\tightlist
\item
  证明 Kolmogoroff 空间（相应地，可达空间）的每个子空间是 Kolmogoroff 空间（相应地，可达空间）；并且比 Kolmogoroff 空间（相应地，可达空间）的某个拓扑更细的拓扑，仍是 Kolmogoroff 空间（相应地，可达空间）的拓扑。
\end{enumerate}

¶ 5) a) 设 X 是无限集。证明 X 上所有拓扑 \(\mathcal{G}_{\mathfrak{c}}\)（§ 2，习题 8），只要 \(1 < c \leqslant\) card (X)，就都是可达空间的拓扑，但都不是豪斯多夫的。

\begin{enumerate}
\def\labelenumi{\alph{enumi})}
\setcounter{enumi}{1}
\tightlist
\item
  证明 X 上所有豪斯多夫拓扑之交是非豪斯多夫拓扑 \(\mathfrak{G}_{\mathrm{card}(\mathbf{N})}\)，它也是 X 上最粗的可达空间拓扑。（利用§ 2 的习题 8，并注意 X 上存在这样的豪斯多夫拓扑：关于它，X 有非闭的可数无限子集。）
\end{enumerate}

\begin{enumerate}
\def\labelenumi{\arabic{enumi})}
\setcounter{enumi}{5}
\tightlist
\item
  \begin{enumerate}
  \def\labelenumii{\alph{enumii})}
  \tightlist
  \item
    设 X 是豪斯多夫空间，D 是 X 的稠密子集。证明 \(\text{card}(X) \leqslant 2^{2^{\text{card}(D)}}\)（注意 X 的每一点都是 D 上某个滤子基的极限）。证明 card (X) 的这个上界不能改进（§ 4，习题 5 b)），并由此推出：在无限集 A 上，超滤子所成的集与滤子所成的集都和 \(\mathfrak{P}(\mathfrak{P}(A))\) 等势。
  \end{enumerate}
\end{enumerate}

\begin{enumerate}
\def\labelenumi{\alph{enumi})}
\setcounter{enumi}{1}
\tightlist
\item
  设 K 是由数 o 和 1 组成的离散空间，A 是无限集。由 a) 推出：若 card (A) ≥ 2 \(^{2\text{card(N)}}\)，则积空间 K \(^{\text{A}}\) 满足§ 1 习题 7 的条件 (D \(_{\text{IV}}\) )，但既不满足 (D \(_{\text{II}}\) ) 也不满足 (D \(_{\text{III}}\) )。（§ 4，习题 4。）
\end{enumerate}

* 7) 设 X 是 R 中的区间 \([-1, 1]\)。考虑 X 上如下的等价关系 S：若 \(x \neq \pm 1\)，则 \(x\) 的等价类由 \(x\) 和 \(-x\) 组成；\(1\)（相应地，\(-1\)）的等价类仅由 \(1\)（相应地，\(-1\)）组成。证明 S 是开的，且商空间 X/S 可达但非豪斯多夫。再证明 X/S 的每一点在 X/S 中都有一个作为正则子空间的邻域。*

\begin{enumerate}
\def\labelenumi{\arabic{enumi})}
\setcounter{enumi}{7}
\tightlist
\item
  \begin{enumerate}
  \def\labelenumii{\alph{enumii})}
  \tightlist
  \item
    设 X 是子空间族 \((\mathbf{X}_{t})\) 的和所成的豪斯多夫（相应地，正则）空间，X/R 是把诸 \(X_{t}\) 沿闭子集 \(A_{tx}\) 借助同胚 \(h_{x_{t}}\)（§ 3，第 4 目）粘合起来而得的商空间；再设对每个指标 t，使得 \(A_{tx} \neq \emptyset\) 的指标 x 只有有限个。在这些条件下，证明 X/R 是豪斯多夫（相应地，正则）的。
  \end{enumerate}
\end{enumerate}

* b) 证明习题 7 的空间 X/S 可以通过沿两个开集把两个正则子空间粘合在一起而得到。*

\begin{enumerate}
\def\labelenumi{\arabic{enumi})}
\setcounter{enumi}{8}
\tightlist
\item
  设 \(\Gamma\) 是豪斯多夫（相应地，正则）空间 X 的一个有限同胚群，R 是由“存在 \(\sigma \in \Gamma\) 使得 \(y = \sigma(x)\)”（X 的点 \(x\) 与 \(y\) 之间）所定义的等价关系。证明商空间 X/R 是豪斯多夫（相应地，正则）的。
\end{enumerate}

* 10) 设 X 是 R 的这样的子空间：它是点 \(\mathbf{I} / n\) 所成之集的余集，其中 \(n\) 是异于 \(\mathbf{o}, \mathbf{i}\) 和 \(-\mathbf{i}\) 的有理整数。设 S 是 X 上的等价关系，其等价类为 (i) 集合 \(\{\mathbf{o}\}\)；(ii) 全体非零整数所成之集；(iii) 所有形如 \(\{x, \mathbf{i} / x\}\) 的集合，其中 \(x \in \mathbf{X}\) 且 \(\mathbf{o} < |x| < \mathbf{i}\)。

\begin{enumerate}
\def\labelenumi{\alph{enumi})}
\item
  证明 X 上的关系 S 既非开的也非闭的。
\item
  证明 S 的图像在 \(X \times X\) 中闭，但商空间 X/S 不是豪斯多夫的。*
\end{enumerate}

¶ * 11) 设 X 是拓扑空间。用 Z 表示积空间 \(R^{X}\)，并对每个 \(x \in X\)，设 \(Z_{x}\) 为 Z 中由所有满足如下条件的 \(u \in R^{X}\) 组成的子集：\(u(x) \in Q\)，且 \(u(y) \notin Q\) 对所有异于 x 的 \(y \in X\) 成立。设 Y 是积 \(X \times Z\) 的子空间，它是诸集 \(\{x\} \times Z_x\)（当 \(x\) 取遍 X 时）的并。证明 Y 是豪斯多夫的，且 X 同胚于 Y 的一个商空间。*

\begin{enumerate}
\def\labelenumi{\arabic{enumi})}
\setcounter{enumi}{11}
\tightlist
\item
  设 X 是至少有两点的拓扑空间。
\end{enumerate}

\begin{enumerate}
\def\labelenumi{\alph{enumi})}
\item
  证明拓扑 \(\mathfrak{C}_{\Omega}\) 和 \(\mathfrak{C}_{\Phi}\)（在 \(\mathfrak{P}_0(\mathbf{X})\) 上；§ 2，习题 7）都不是可达空间的拓扑。
\item
  设 \(\mathcal{G}_{\Theta}\) 表示拓扑 \(\mathcal{G}_{\Omega}\) 与 \(\mathcal{G}_{\Phi}\) 的上确界。证明集 \(\mathfrak{F}(\mathbf{X})\)（\(\mathbf{X}\) 的非空闭子集所成之集）赋以由 \(\mathcal{G}_{\Theta}\) 诱导的拓扑后是 Kolmogoroff 空间；且若 \(\mathbf{X}\) 可达，则 \(\mathfrak{F}(\mathbf{X})\) 也可达。
\item
  在 X 可达的假设下，证明空间 \(\mathfrak{F}(\mathbf{X})\)（赋以由 \(C_{\Theta}\) 诱导的拓扑）是豪斯多夫的当且仅当 X 正则。
\end{enumerate}

¶ 13) 设 X 是一个集合，Γ 是 X 的一个置换群。a) 设 F 是 X 的子集所成的一个集。设 S 是 X 上使 F 中所有集合皆闭且每个 u ∈ Γ 都是同胚的所有拓扑组成的集。证明集 S 有一个最小元 C(Γ, F)。若 X 无限，证明存在 X 的子集所成的集 F，使得 C(Γ, F) 可达但非离散（§ 2，习题 8）。

\begin{enumerate}
\def\labelenumi{\alph{enumi})}
\setcounter{enumi}{1}
\tightlist
\item
  设 \(J(\Gamma)\) 是具有下述性质的子集 \(A\)（属于 \(X\)）所成的集：对每个 \(x \notin A\)，存在映射 \(f \in \Gamma\)，使得 \(f(x) \neq x\) 且保持 \(A\) 的每个元素不动。证明 \(X\) 上使每个映射 \(u \in \Gamma\) 都是 \(X\) 到自身的同胚的每个豪斯多夫拓扑，都比 \(\mathcal{C}_{J}(\Gamma) = \mathcal{C}(\Gamma, J(\Gamma))\) 更细。
\end{enumerate}

证明拓扑 \(\mathfrak{G}_{\mathrm{J}}(\Gamma)\) 是豪斯多夫的当且仅当群 \(\Gamma\) 满足下列条件：对任意两个不同的元素 \(x, y\)（属于 \(\mathbf{X}\)），存在有限个元素 \(u_{k} \in \Gamma\)，使得（记 \(\mathbf{F}(u_{k})\) 为 \(\mathbf{X}\) 中在 \(u_{k}\) 下不变的元素所成之集）诸 \(\mathbf{F}(u_{k})\) 覆盖 \(\mathbf{X}\)，且没有一个 \(\mathbf{F}(u_{k})\) 同时含 \(x\) 和 \(y\)。

\begin{enumerate}
\def\labelenumi{\alph{enumi})}
\setcounter{enumi}{2}
\item
  设 \(\mathcal{C}_{0}\) 表示有理直线 \(\mathbf{Q}\)（*相应地，实直线 \(\mathbf{R}_{*}\)*）的拓扑，并设 \(\Gamma\) 是 \(\mathbf{Q}\)（*相应地，\(\mathbf{R}_{*}\)*）到自身的所有同胚组成的群。证明拓扑 \(\mathcal{C}_{\mathbf{J}}(\Gamma)\) 与 \(\mathcal{C}_{0}\) 相同。
\item
  设 X 是有理直线的由点 0 和诸点 \(\iota/n\)（n 为整数 \(\geqslant\iota\)）组成的子空间，并设 \(\Gamma\) 是 X 的同胚群。证明拓扑 \(\mathcal{G}_{\mathrm{J}}(\Gamma)\) 不是豪斯多夫的，且 X 关于这个拓扑的同胚群就是 \(\Gamma\)。
\item
  设 X 是这样的拓扑空间：对 X 的每对不同的点 x, y，存在 X 到自身的同胚 u，使得 \(u(x) = x\) 且 \(u(y) \neq y\)（*例如空间 \(R^{n}_{*}\)*）。设 G 是 X 到自身的同胚群，并对每个 \(x \in X\)，设 \(S_{x}\) 是 G 中保持 x 不动的子群。证明 \(S_{x}\) 等于它在 G 中的正规化子，且所有子群
\end{enumerate}

\(S_{x}\) 之交仅由单位元组成；特别地，G 的中心仅由单位元组成。

\begin{enumerate}
\def\labelenumi{\arabic{enumi})}
\setcounter{enumi}{13}
\tightlist
\item
  证明公理 \((\mathrm{O}_{\mathrm{III}})\) 与下列每条公理等价：
\end{enumerate}

\((\mathrm{O}_{\mathrm{III}}^{\prime\prime})\) 若 F 是 X 的任一闭子集，则 F 的所有闭邻域之交等于 F。

\((\mathrm{O}_{\mathrm{III}}^{\prime\prime})\) 若 F 是 X 上收敛于点 a 的任一滤子，则以 F 中诸集的闭包为基的 X 上的滤子 F 也收敛于 a。

\begin{enumerate}
\def\labelenumi{\arabic{enumi})}
\setcounter{enumi}{14}
\tightlist
\item
  \begin{enumerate}
  \def\labelenumii{\alph{enumii})}
  \tightlist
  \item
    证明满足公理 (O \(_{III}\) ) 的每个 Kolmogoroff 空间都是豪斯多夫的（从而是正则的）。
  \end{enumerate}
\end{enumerate}

\begin{enumerate}
\def\labelenumi{\alph{enumi})}
\setcounter{enumi}{1}
\tightlist
\item
  在一个三元集上构造一个满足公理 \((\mathrm{O}_{\mathrm{III}})\) 的非豪斯多夫拓扑。
\end{enumerate}

\begin{enumerate}
\def\labelenumi{\arabic{enumi})}
\setcounter{enumi}{15}
\item
  设 \((\mathbf{X}_t)\) 是拓扑空间的一个无限族，并在积集合 \(\mathbf{X} = \prod_{i} \mathbf{X}_i\) 上赋以§ 4 习题 9 中定义的拓扑之一。证明若诸 \(\mathbf{X}_t\) 是 Kolmogoroff（相应地，可达、豪斯多夫、正则）空间，则 \(\mathbf{X}\) 亦然。
\item
  证明拓扑 \(\mathcal{C}_0(\mathbf{X})\)、\(\mathcal{C}_+(\mathbf{X})\) 和 \(\mathcal{C}_{-}(\mathbf{X})\)（§ 2，习题 5）在线性序集 \(\mathbf{X}\) 上都是正则的。
\item
  设 \(X_{1}\)，\(X_{2}\) 是两个集合，\(F_{i}\) 是 \(X_{i}\) 上的滤子（i = 1, 2），f 是 \(X_{1} \times X_{2}\) 到正则空间 Y 中的映射，使得对每个
\end{enumerate}

\[
x _ {1} \in \mathrm{X} _ {1} \text { there   exists } \lim _ {x _ {2}, \mathfrak {F} _ {2}} f (x _ {1}, x _ {2}) = g (x _ {1}).
\]

证明：点 \(a \in Y\) 是 \(g\) 关于滤子 \(\mathfrak{F}_1\) 的聚点，当且仅当任给一个邻域 \(V\)（\(a\) 在 \(Y\) 中的）和任一集 \(\mathbf{A}_1 \in \mathfrak{F}_1\)，存在 \(x_1 \in \mathbf{A}_1\) 及 \(\mathbf{A}_2 \in \mathfrak{F}_2\)，使得 \(f(\{x_1\} \times \mathbf{A}_2) \subset V\)。¶ 19) 设 \(X\) 是非正则的豪斯多夫空间（参见习题 20），\(a\) 是 \(X\) 的一点，使得存在邻域 \(U\)（\(a\) 的），它不含 \(a\) 的任何闭邻域。设 \(\mathfrak{F}\) 是 \(a\) 的邻域滤子，\(Y = \mathfrak{F} \cup \{\omega\}\) 是与 \(\mathfrak{F}\) 的截面滤子相伴随的拓扑空间。设 \(Z\) 是集合 \(Y \times X\)，赋以如下拓扑：对任何点 \((y, z) \neq (\omega, a)\)，\((y, z)\) 的邻域与积拓扑下的相同，而 \((\omega, a)\) 的邻域是含有形如

\[
(\{\omega \} \times \mathrm{V}) \cup ((\mathrm{Y} - \{\omega \}) \times \mathrm{X}),
\]

之集，其中 \(\mathbf{V}\) 是 \(a\) 在 \(\mathbf{X}\) 中的邻域。设 \(\mathbf{A}\) 是 \(\mathbf{Z}\) 的子集，它是诸集 \(\{\mathbf{V}\} \times \mathbf{V}\)（\(\mathbf{V} \in \mathfrak{F}\)）的并。设 \(f\) 是 A 到 X 中的映射 \((V, x) \rightarrow x\)。证明 f 在 \(\overline{A}\) 的每一点处关于 A 都有极限，但不存在 \(\overline{A}\) 到 X 中的连续映射延拓 f。

¶ 20) a) 拓扑空间 X 的开子集 U 称为正则的，如果 \(\mathbf{U} = \alpha(u)\)（§ 1，习题 3）；等价地，U 是正则的开集，如果 U 是某个闭集的内部。如果 X 的正则开子集构成 C 的一个基，则称 X 是半正则的（并称它的拓扑 C 是半正则的）。证明：若 C 满足公理 \((\mathrm{O}_{\mathrm{III}})\)，则 C 是半正则的。

\begin{enumerate}
\def\labelenumi{\alph{enumi})}
\setcounter{enumi}{1}
\item
  证明两个正则开集的交是正则开集。由此推出：关于 \(\mathcal{C}\) 的正则开集构成 X 上一个拓扑 \(\mathcal{C}^*\) 的基，该拓扑比 \(\mathcal{C}\) 粗且半正则；\(\mathcal{C}^*\) 称为与 \(\mathcal{C}\) 相伴随的半正则拓扑。证明：对 X 的每个在 \(\mathcal{C}\) 中开的子集 U，U 关于 \(\mathcal{C}^*\) 的闭包与它在 \(\mathcal{C}\) 中的闭包相同；且若 X 关于 \(\mathcal{C}\) 可达，则 X 关于 \(\mathcal{C}^*\) 的孤立点与关于 \(\mathcal{C}\) 的孤立点相同。证明 \(\mathcal{C}^*\) 是豪斯多夫的当且仅当 \(\mathcal{C}\) 是豪斯多夫的{[}利用§ 1 的习题 3 d){]}；最后证明：若 Y 是正则空间，则 X 到 Y 中的连续映射在拓扑 \(\mathcal{C}\) 与 \(\mathcal{C}^*\) 下是相同的。
\item
  设 X 是半正则空间，\(C_{0}\) 是它的拓扑。证明 X 上的拓扑 C 满足 \(C^{*} = C_{0}\)，当且仅当存在 X 的稠密子集（在拓扑 \(C_{0}\) 下）所成的一个集 M，使得 M 中诸集的每个有限交属于 M，且拓扑 C 由 M 与在拓扑 \(C_{0}\) 下为开的 X 的子集所成之集的并生成。{[}考虑（拓扑 C 下的）稠密开子集，并注意 C 中的每个开集是一个（拓扑 C 下的）稠密开集与一个（拓扑 \(C_{0}\) 下的）开集的交。{]}由此构造比正则拓扑更细的非正则豪斯多夫拓扑的例子（*甚至比可度量化紧空间的拓扑更细*）。
\end{enumerate}

\begin{enumerate}
\def\labelenumi{\arabic{enumi})}
\setcounter{enumi}{20}
\tightlist
\item
  设 X 是没有孤立点的豪斯多夫空间，\(C_{0}\) 是 X 的拓扑。设 M 表示 X 的这样的可数子集 A 的余集所成之集：\(\overline{A}\) 只有有限个非孤立点。证明：若 C 是 X 上由 M 与在 \(C_{0}\) 下为开的子集所成之集的并生成的拓扑，则 \(C^{*} = C_{0}\)（习题 20），且关于 C 收敛的每个序列只有有限个不同的项（这蕴含关于 C，X 的任何点都没有可数的基本邻域系，尽管 X 本身可以是可数的，从而 X 的每一点都是该点的可数个邻域之交）。
\end{enumerate}

¶ 22) a) 用习题 20 c) 的记号，证明拓扑的集 \(\mathrm{E}(\mathfrak{C}_{0})\)（由满足 \(C^{*}=C_{0}\) 的拓扑 C 组成）关于关系“C 比 \(C'\) 粗”是归纳的。若 \(C_{0}\) 是 X 上的半正则拓扑，则 \(\mathrm{E}(\mathfrak{C}_{0})\) 的一个极大元称为次极大拓扑，而赋以这种拓扑的集合 X 称为次极大空间。于是 \(\mathrm{E}(\mathfrak{C}_{0})\) 中的每个拓扑都比某个次极大拓扑粗。

\begin{enumerate}
\def\labelenumi{\alph{enumi})}
\setcounter{enumi}{1}
\item
  \(\mathfrak{C}\) 是 \(\mathbf{X}\) 上的次极大拓扑，当且仅当 \(\mathbf{X}\) 的每个在拓扑 \(\mathfrak{C}\) 下稠密的子集在 \(\mathfrak{C}\) 中都是开的。因此非空的次极大空间不是可解的（§ 2，习题 2）。
\item
  证明次极大空间的每个子空间是局部闭且次极大的（先考虑开子空间与闭子空间）。
\item
  证明与一个滤子相伴随的空间是次极大空间，这里该滤子比无限集的有限子集之余集所成的滤子更细。
\item
  若 M 是次极大空间 X 的子集，证明子空间 \(\overline{M} - \mathring{M} = \operatorname{Fr}(M)\) 是离散的。
\item
  反之，设 X 是这样的拓扑空间：它有一个稠密开子集 U，使得 U 是次极大的，且 X 的拓扑是 U-极大的（§ 3，习题 11）。那么 X 是次极大的。
\end{enumerate}

\begin{enumerate}
\def\labelenumi{\arabic{enumi})}
\setcounter{enumi}{22}
\item
  设 X 是拓扑空间，A 是 X 的稠密子集，使得 X 的拓扑是 A-极大的（§ 3，习题 11）。设 f 是 A 到拓扑空间 Y 中的连续映射，使得在 X --- A 的每一点处，f 关于 A 的聚点集非空（*例如当 X 非空且 Y 拟紧时即是这种情形*）。证明 f 可以连续延拓到整个 X 上。
\item
  在有理直线的区间 {]}0,1{[} 中，设 A 是形如 \(k/2^n\) 的有理数所成之集，B 是形如 \(k/3^n\) 的有理数所成之集（\(k\) 为整数）。设 X 是在 {]}0,1{[} 上赋以如下拓扑而得的空间：该拓扑由包含在 X 中的开区间 {]}α,β{[} 以及集合 A 和 B 生成。证明 X 是豪斯多夫且半正则的（习题 20），但不是正则的。
\end{enumerate}

¶ 25 a) 在集合 X 上，正则拓扑所成之集与没有孤立点的正则拓扑所成之集关于关系“C 比 \(C'\) 粗”都是归纳集。

\begin{enumerate}
\def\labelenumi{\alph{enumi})}
\setcounter{enumi}{1}
\item
  集合 X 上的拓扑 \(\mathcal{G}\) 称为超正则的，如果它在 X 上正则且无孤立点的拓扑所成的集中是极大的；因此对 X 上任何无孤立点的正则拓扑，都存在比它更细的超正则拓扑。空间称为超正则的，如果它的拓扑是超正则的。证明无孤立点的正则拓扑 \(\mathcal{G}\) 是超正则的，当且仅当只要 A 与 B 是 X 的两个互补的无孤立点子集，A 和 B 就都是开集。特别地，超正则空间不是可解的（§ 2，习题 11）。（为证条件是充分的，注意：若 \(\mathfrak{C}'\) 是比 \(\mathfrak{G}\) 更细的无孤立点正则拓扑，U 是关于 \(\mathfrak{C}'\) 的非空开集，\(\overline{\mathrm{U}}\) 是它在拓扑 \(\mathfrak{C}'\) 下的闭包，则 \(\overline{\mathrm{U}}\) 与 \(\mathrm{X} - \overline{\mathrm{U}}\) 关于拓扑 \(\mathfrak{C}'\) 都没有孤立点。）
\item
  在超正则空间中，每个无孤立点的集合的闭包是开的，且稠密集的内部稠密；因此两个稠密集之交稠密。由此推出：若 \(C_{0}\) 是超正则拓扑，则在所有满足 \(C^{*} = C_{0}\)（习题 20）的拓扑 C 之中，有一个（必然是次极大的）比所有其他拓扑都细。
\end{enumerate}

* 26) 设 \((\theta_{n})\) 是无理数序列，在区间 \(\mathbf{I} = [\mathbf{o},\mathbf{i}]\)（\(\mathbf{R}\) 中的）内稠密。设 \(\mathbf{I}_n\) 是把 \(\mathbf{I}\) 的所有点 \(0, \theta_1, \ldots, \theta_n\) 等同起来而得的商空间，\(\varphi_n\) 是典范映射 \(\mathbf{I} \to \mathbf{I}_n\)。设 \(\mathbf{X}\) 是含于 \(\mathbf{I}\) 的有理数所成之集；则 \(\varphi_n\) 在 \(\mathbf{X}\) 上的限制是 \(\mathbf{X}\) 到 \(\varphi_n(\mathbf{X})\) 上的双射。设 \(\mathcal{G}_n\) 是在 \(\varphi_n(\mathbf{X})\) 上由 \(\mathbf{I}_n\) 的拓扑诱导的拓扑经该双射的逆像。证明递减序列 \((\mathcal{G}_n)\)（\(\mathbf{X}\) 上的豪斯多夫拓扑序列）的下确界不是豪斯多夫拓扑。*

¶ 27) 设 X 是拓扑空间。证明存在 Kolmogoroff（相应地，可达、豪斯多夫、正则）空间 Y 和连续映射 \(\varphi: X \to Y\)，它具有下列泛性质：对 X 到 Kolmogoroff（相应地，可达、豪斯多夫、正则）空间 Z 中的每个连续映射 \(f: X \to Z\)，存在唯一的连续映射 \(g: Y \to Z\)，使得 \(f = g \circ \varphi\)（《集合论》，第四章，§ 3，第 1 目）。证明：

\begin{enumerate}
\def\labelenumi{\alph{enumi})}
\item
  对 Kolmogoroff 空间，可取 Y 为 X 关于等价关系 \(\{\bar{x}\}=\{\bar{y}\}\) 的商空间。
\item
  对可达空间，可取 Y 为 X 关于下述等价关系 R 的商空间：R 的图像是 X 上所有满足 \(\mathrm{G}(x)\) 对每个 \(x \in X\) 在 x 中闭的等价关系之图像 G 的交；于是 R 的图像含有 x 与 y 之间的关系 \(\{\bar{x}\} \cap \{\bar{y}\} \neq \emptyset\) 的图像。
\item
  对豪斯多夫空间与正则空间，试证《集合论》第四章，§ 3，第 2 目的条件 \((\mathrm{CU}_{\mathbf{I}})\) 至 \((\mathrm{CU}_{\mathbf{III}})\) 都满足（特别要利用习题 6 a)）。给出这样的例子：X 不正则，而 X 到相应的泛正则空间中的映射 \(\varphi: X \to Y\) 是双射{[}参见习题 20 b){]}。
\end{enumerate}

\begin{enumerate}
\def\labelenumi{\alph{enumi})}
\setcounter{enumi}{3}
\tightlist
\item
  若 X 是满足公理 \((\mathrm{O}_{\mathrm{III}})\) 的拓扑空间，证明与 X 对应的泛 Kolmogoroff 空间典范同胚于与 X 对应的泛正则空间。
\end{enumerate}

\begin{enumerate}
\def\labelenumi{\arabic{enumi})}
\setcounter{enumi}{27}
\tightlist
\item
  设 X 是非正则的豪斯多夫空间，F 是 X 的闭子集，\(a \notin F\) 是 X 的一点，使得 a 的每个邻域都与 F 的每个邻域相交。设 R 是把 F 的所有点等同起来而得的 X 上的等价关系。证明 R 是闭的，且 R 的图像在 \(X \times X\) 中闭，但 R 不是豪斯多夫的。
\end{enumerate}

¶ 29) a) 设 R 是非空拓扑空间 X 上的豪斯多夫等价关系，\(\mathfrak{F}(X)\) 是 X 的非空闭子集所成的空间，赋以由 \(\mathfrak{G}_{\Theta}\)（习题 12）诱导的拓扑。证明每个位于 X/R 的闭包中之 \(A \in \mathfrak{F}(X)\)（{[}把 X/R 看作 \(\mathfrak{F}(X)\) 的子集{]}）都含于 R 的某个等价类中（用反证法）。

\begin{enumerate}
\def\labelenumi{\alph{enumi})}
\setcounter{enumi}{1}
\tightlist
\item
  再设 X 正则且 R 是开的。证明 X/R 于是是 \(\mathfrak{F}(X)\) 关于拓扑 \(\mathfrak{G}_{\Theta}\) 的闭子集（利用§ 7 第 6 目的命题 11）。
\end{enumerate}

